% Lecture visuelle directe avec la vision du modèle ; aucun OCR.
% Source : PDF original, vue 046, côté gauche, page manuscrite 92, après le chapitre premier.
% Statut : transcrit ; relecture indépendante à effectuer.
% Chapitre complet : p. 92 (fin)–111 (début).
\chapter{Préfecture de Monsieur Ripoud}
\begin{center}
Établissement des sections ; célébration de la fête patronale ; fondation de plusieurs Congrégations, \&c.
\end{center}

\section*{Préfecture de M. Ripoud}
% Mention marginale de la page manuscrite 92.
La Congrégation existait, mais c'était encore bien peu de chose. Elle n'eût pas été capable de donner aucun signe de vie. Elle ressemblait pour lors à un foible ruisseau qui vient à peine de sortir de sa source et qui se complaît
% Source : PDF original, vue 046, côté droit, page manuscrite 93.
à promener lentement son onde tranquille au milieu d'une prairie émaillée de fleurs. C'était un enfant au maillot, incapable de se soutenir par lui-même. Sa bonne mère, la Ste Vierge, versait abondamment sur lui le lait des consolations spirituelles ; les assemblées générales, les fêtes, les communions générales, les réceptions des nouveaux frères, la tendre amitié qui unissait tous les congréganistes répandaient sur leur existence un charme d'autant mieux senti que jusqu'alors il leur avait été inconnu. Le P. de la Congrégation leur adressait fréquemment des instructions pathétiques qui pénétraient leurs âmes de sentimens de la dévotion la plus tendre. Ce bon père était vraiment ce serviteur fidèle que Dieu a préposé à la garde de sa maison ; Marie l'avait chargé de veiller sur la Congrégation naissante ; il était devenu pour elle comme une nourrice attentive qui se regarde comme la seconde mère de l'enfant qui lui est confié ; il connaissait le secret d'apprendre à la jeunesse à servir le Seigneur avec une grande liberté d'esprit, une confiance toute filiale et une joie toute céleste : c'était bien là embellir les premiers jours de la Congrégation par l'aimable sourire de l'enfance. \textit{Gaudete in Domino semper}, nous disait-il avec l'apôtre : \textit{Iterum dico, gaudete}\footnote{Traduction : « Réjouissez-vous toujours dans le Seigneur ; je le dis encore, réjouissez-vous. » Référence : Philippiens 4, 4. (Note du transcripteur).}.

Toutefois malgré la suavité de ce premier état la Congrégation n'était pas seulement destinée à en savourer les douceurs. Elle n'avait pas été formée pour demeurer constamment comme Marie dans la contemplation aux pieds du Sauveur. Il fallait aussi que comme Marthe elle se préparât à le servir d'une manière active. Dieu pour fortifier cet athlète et donner de la souplesse à tous ses mouvements avait versé
% Source : PDF original, vue 047, côté gauche, page manuscrite 94.
sur lui l'huile bienheureuse des béatitudes du cœur. Mais le moment était venu où il devait se revêtir de ses armes et entrer courageusement dans l'arène afin de combattre avec vigueur ces bêtes sauvages, ces sangliers furieux, ces loups dévorants qui ne cessent de ravager l'héritage du Seigneur. Il fallait que l'enfant apprît à marcher, qu'il essayât d'abord ses forces en tremblant et qu'enfin il devînt homme, qu'il s'élançât dans la carrière prête à s'ouvrir devant lui.

\section*{Établissement des sections}
% Mention marginale de la page manuscrite 94, après celle sur la préfecture.
C'est donc afin que la Congrégation commençât à accomplir les desseins qu'il avait formés sur elle que Dieu lui inspira d'organiser dans son sein un service régulier de bonnes œuvres. En se livrant à ces œuvres dans l'intention d'étendre le plus qu'ils le pourraient le règne du Seigneur, les congréganistes remplissaient complètement le but qu'ils s'étaient proposé ; ils procuraient la gloire de Dieu, ils rendaient honneur à Marie sous la protection de laquelle ces œuvres devaient être placées, et enfin ils travaillaient à leur sanctification, car il est écrit : celui qui aura sauvé l'âme de son frère sauvera également la sienne ; ainsi tout en se tenant prête à profiter des circonstances extraordinaires qui pourraient se présenter de travailler à la gloire de Dieu d'une manière plus directe, la Congrégation devait trouver dans cette organisation une occasion de tous les jours de remplir sa mission. De plus l'exercice journalier de ces bonnes œuvres allait amener des réunions fréquentes entre les congréganistes et en donnant à leur zèle un aliment quotidien renouveller sans cesse les forces vitales de la Congrégation.

La divine sagesse inspira donc aux chefs de la Congrégation
% Source : PDF original, vue 047, côté droit, page manuscrite 95.
la pensée salutaire de la diviser en quatre sections chargées chacune d'une bonne œuvre spéciale. Ils proposèrent cette grande mesure à l'assemblée générale de la Congrégation qui s'empressa de l'adopter.

La première section fut appelée la section du zèle. En l'établissant on songea à remplir une lacune qui existait dans le règlement ; il chargeait les deux assistants de veiller aux réceptions ; cela fut trouvé insuffisant ; on confia donc ce soin à cette première section ; on ajouta celui de travailler à la conversion de la jeunesse. La seconde section fut celle de l'instruction ; elle devait par tous les moyens possibles procurer la connaissance salutaire des vérités de l'Évangile à ceux qui ont le malheur de les ignorer. La troisième section, sous le nom de section des aumônes, fut chargée de distribuer des secours aux pauvres honteux ; enfin la quatrième nommée section des prisons et hôpitaux devait exercer les mêmes œuvres de charité à l'égard des prisonniers et des malades de l'hôtel-Dieu. Chaque section devait avoir un chef à la nomination du préfet. Tous les membres sans autre exception que le préfet seul devaient être classés dans une des quatre sections. Vous voyez, Messieurs, que la manière dont cet arrêté a établi les sections diffère un peu de la législation actuellement en vigueur. Nous aurons soin de signaler plus tard les modifications qui ont été jugées convenables et les motifs ou les circonstances qui les ont occasionnées.

Cet arrêté qui donnait une physionomie toute nouvelle à la Congrégation reçut de suite son exécution. Mais vous devez bien penser que nos pauvres petites sections ne commencèrent pas leurs travaux d'une manière bien active. La Congrégation
% Source : PDF original, vue 048, côté gauche, page manuscrite 96.
était animée d'un grand zèle, mais elle manquait d'expérience. Il lui était souvent difficile de se procurer les moyens d'exécution ; elle tâtonnait un peu et parfois faisait des pas d'école ; comme je vous le disais tout à l'heure elle était encore dans l'enfance. Elle vacillait à chaque pas ; qu'en serait-il advenu si la bonne mère ne fût venue à son secours ? Mais la main de Marie était là et nous verrons bientôt que malgré les hésitations des premiers temps toutes les difficultés finirent par être vaincues. Cependant la section de l'instruction avait commencé à faire le catéchisme à un jeune homme ; celle des aumônes visitait et portait quelques foibles secours à un ou deux pauvres ; celle des prisons et hôpitaux rendait quelques visites à deux ou trois malades de l'hôtel-Dieu ; elle n'avait pas pu parvenir à se faire ouvrir encore les portes des prisons, mais en attendant elle conservait précieusement dans sa caisse son trésor montant à 4~f.~50 qu'elle économisait pour l'hiver. J'entre avec simplicité dans ces menus détails, Messieurs, pour que vous soyez bien convaincus que dans la Congrégation comme lorsqu'il a créé le monde tout ce que Dieu a fait il l'a fait de rien. Nous reviendrons plus tard sur ces trois sections pour admirer dans leurs progrès la marche de la Providence.

La section du zèle était sans contredit en possession de l'œuvre à laquelle la Congrégation devait porter alors le plus d'intérêt, aussi ce fut elle qui déploya de suite le plus d'activité. À chaque nouvelle fête une réception venait augmenter le nombre des frères ; pour préparer d'avance ces réceptions et parvenir à mieux connaître ceux dont elle s'occupait et
% Source : PDF original, vue 048, côté droit, page manuscrite 97.
en même temps pour essayer suivant le vœu de l'arrêté qui l'avait constituée de ramener à la religion les jeunes gens, elle imagina de former une espèce de cercle où elle pût recevoir les étrangers. Cela se fit sans aucuns frais ; un de nos amis prêta son salon ; on s'y réunissait trois fois par semaine ; mais cette tentative comme presque toutes celles du même genre qui ont été essayées depuis resta complètement sans succès ; on ne tarda pas à y renoncer.

\section*{Fête de la Conception immaculée}
% Mention marginale de la page manuscrite 97.
L'organisation des sections et les foibles commencements de leurs travaux ne nous faisaient pas oublier nos solennités. Nous y trouvions trop de jouissances pour ne pas éprouver constamment un désir bien vif de les voir revenir. Elles tardaient toujours trop à notre impatience. Plusieurs jours d'avance nos imaginations en étaient fortement occupées ; la veille chacun de nous venait à l'envi travailler à la décoration de l'autel au pied duquel nous devions nous réunir. Tous ceux des congréganistes de 1802 qui restent encore dans vos rangs n'ont pas oublié plus que moi, je pense, cette joie si douce à nos cœurs avec laquelle nous vîmes arriver l'époque dans laquelle pour la première fois nous devions célébrer notre fête patronale, la fête de la Conception immaculée de la Ste Vierge. Nous étions si heureux de rendre tous ensemble cet hommage solennel à ce glorieux privilège de notre mère ! Tout fut disposé pour cette fête dans le salon d'un de nos amis. Ce salon était par lui-même assez élégamment décoré ; on en condamna les croisées pour intercepter la lumière du jour ; on le drapa avec des rideaux de soie cramoisie ; des lustres placés de distance en distance répandaient une religieuse clarté ; au fond de la chapelle s'élevait un autel resplendissant de lumières et garni
% Source : PDF original, vue 049, côté gauche, page manuscrite 98.
d'une multitude de vases de fleurs ; tout autour des guirlandes de fleurs se dessinaient en élégants festons. Au milieu de cette gracieuse décoration on apercevait ces mots tracés en lettres d'or : \textit{Maria semper immaculata}\footnote{Traduction : « Marie toujours immaculée. » (Note du transcripteur).}. La Congrégation avait voulu offrir aussi à sa mère comme son bouquet pour le jour de sa fête le renouvellement de la déclaration solennelle par laquelle l'Église de Lyon reconnut autrefois en établissant une fête pour le célébrer le glorieux triomphe que par sa Conception immaculée la mère de Dieu a remporté sur l'auteur du péché. Notre société était heureuse de prendre ainsi dans les temps modernes comme l'initiative d'une dévotion qui depuis quelques années est devenue la dévotion de choix et de prédilection de toutes les âmes pieuses.

Pardon, Messieurs, mais qu'il me soit permis de suspendre ici un instant mon récit. Pourrais-je rappeler le souvenir de cette fête de la Conception célébrée pour la première fois par la Congrégation sans élever avec vous mon âme vers notre Dieu en ma qualité de Lyonnais dans un profond sentiment de reconnaissance ? Oh ! de quels sentimens ne devons-nous pas être animés lorsque nous réfléchissons sur la grâce particulière dont notre ville a été l'objet. Il a plu au Seigneur de la placer au nombre des villes privilégiées chères à Marie à qui il a été donné de connaître les premières et de révéler au monde chrétien les plus glorieuses prérogatives de la mère de Dieu ! Qu'il nous soit permis de croire que c'est à l'heureuse intervention de St Jean l'Évangéliste, notre père dans la foi puisqu'il nous a envoyé nos premiers apôtres, que nous avons dû ce précieux avantage. Ce disciple bien-aimé auquel la Ste
% Source : PDF original, vue 049, côté droit, page manuscrite 99.
Vierge a été confiée sur le Calvaire comme un précieux dépôt nous a transmis comme autrefois à Éphèse une tendresse filiale envers Marie plus répandue qu'elle ne l'est ailleurs. Oh ! Marie, oh ! notre bonne mère ! Le cœur du Lyonnais, vos enfants, surabonde de joie lorsqu'ils se souviennent que s'il fut donné dans le 5\up{e} siècle à la ville d'Éphèse de s'écrier avec transport : Marie est mère de Dieu, nos pères ont eu la mission de répondre dans le 13\up{e} : oui, Marie est mère de Dieu, nous le disons avec bonheur ; mais Marie peut s'honorer aussi d'un autre privilège et nous sommes heureux de le proclamer : Marie a été conçue sans péché. Et l'Église revêtue de l'autorité de Dieu même est venue sanctionner l'une et l'autre vérité. Ainsi donc, ô ma mère, nous, nous qui sommes vos enfants, nous qui mettons toute notre gloire dans la gloire de notre mère, ce sera avec un saint enthousiasme que nous répéterons ces deux vérités proclamées par Éphèse, proclamées par Lyon, proclamées par toute l'Église catholique : Marie est mère de Dieu ! Marie a été conçue sans péché ! Oh ! notre sainte protectrice ! Vous nous permettez bien d'ajouter dans toute l'effusion de notre tendresse filiale : Marie est aussi notre mère. Nous chercherons alors dans votre cœur maternel un refuge assuré ; nous y trouverons la force nécessaire pour résister aux suggestions de l'ennemi de notre salut et pour supporter les misères de la vie. C'est par votre cœur, ô ma mère, que nous arriverons au cœur de votre adorable fils. C'est encore votre cœur qui à notre moment suprême nous ouvrira la porte du ciel, notre patrie. C'est alors que nous pourrons joindre nos foibles voix à celles des anges et des saints
% Source : PDF original, vue 050, côté gauche, page manuscrite 100.
pour bénir et louer à jamais notre Dieu et le remercier de toute la gloire dont il a revêtu celle que l'archange Gabriel salua du nom de pleine de grâces et de toutes les faveurs qu'elle nous a obtenues de sa main paternelle.

Tous comme aujourd'hui, Messieurs, ces sentimens pénétraient nos cœurs, aussi notre fête fut-elle célébrée avec une sainte allégresse. Il y eut réunion le matin et le soir ; la solennité fut terminée par un salut solennel ; on y chanta les litanies de la Ste Vierge à trois voix ; M. Boni les accompagnait avec sa flûte. Comme je cherche à enregistrer autant que mes souvenirs me le permettront le commencement de tous les usages que le temps a consacrés, je dois ajouter ici que ce fut à cette fête qu'on récita pour la première fois avant la messe matines et laudes de l'office de la Ste Vierge. On eut l'intention de donner ainsi le temps d'arriver avant le commencement de la messe à ceux qui se trouveraient en retard sans toutefois déroger à l'excellente habitude de commencer à heures fixes.

\section*{Congrégation des demoiselles}
% Mention marginale de la page manuscrite 100.
Le même jour notre auguste patronne recevait à Lyon un autre hommage solennel. Une nouvelle Congrégation semblable à la nôtre se fondait en ce jour et venait se placer comme nous sous la protection de la mère de Dieu sous le titre de son immaculée Conception. Cette Congrégation était la Congrégation des demoiselles. Remarquons en passant une nouvelle preuve de la bienveillance de Marie envers ceux qui se constituent les défenseurs d'un privilège si cher à son cœur. Nous verrons bientôt cette nouvelle Congrégation et la nôtre, toutes deux sous le titre de la Conception immaculée, résister seules à l'orage qui ne tardera pas à s'élever,
% Source : PDF original, vue 050, côté droit, page manuscrite 101.
tandis que d'autres sociétés du même genre en ont été renversées. La Congrégation des demoiselles avait encore d'autres droits à une assistance toute particulière du ciel ; les personnes éminemment pieuses qui en formèrent le premier noyau avaient déjà dans les premières années de la Révolution été réunies en Congrégation sous la direction de M. Linsolas. Parmi leurs sœurs de cette époque plusieurs avaient défendu la cause de la religion devant les tyrans sanguinaires de la Terreur et avaient trouvé sur l'échafaud la couronne du martyre. Celles qui avaient survécu étaient dignes par tous les sentimens qui remplissaient leurs âmes d'avoir été unies avec ces héroïnes dans le service de notre commun maître. Elles pouvaient bien se regarder comme certaines de trouver dans ces anciennes amies qu'elles étaient en droit de révérer comme des saintes des protectrices qui seraient les plus fermes soutiens de la Congrégation naissante. Aussi la Congrégation des demoiselles s'est-elle constamment soutenue et a toujours prospéré ; elle s'est sévèrement maintenue dans cette modestie et cette réserve qui fait le plus bel apanage des personnes du sexe ; mais elle n'en a pas moins répandu la bonne odeur de J.-C., car nonobstant les soins qu'elle prenait pour se dérober à tous les yeux ses grandes vertus se trahissent malgré elle par la délicieuse odeur de prédestination qu'elles exhalent.

Les œuvres de Dieu doivent toujours s'attendre à voir naître autour d'elles quelques obstacles ou quelques difficultés. Jusqu'à présent tout semblait marcher au gré de nos désirs ; mais voilà qu'un événement qui cependant était prévu quoiqu'on ne le crût pas aussi rapproché vint
% Source : PDF original, vue 051, côté gauche, page manuscrite 102.
menacer d'étouffer la Congrégation dans son berceau, car elle avait encore si peu de forces que si elle n'avait pas été soutenue d'en haut il n'en fallait pas tant pour l'anéantir.

\section*{Vocation ecclésiastique de trois fondateurs}
% Mention marginale de la page manuscrite 102.
Nous avons dit précédemment que depuis quelques mois Messieurs Ripoud, Jordan et Boni cherchaient à connaître leur vocation. Tous trois se décidèrent à peu près en même temps à quitter le monde pour embrasser l'état ecclésiastique. Or comme vous avez pu le remarquer non seulement tous les trois étaient fondateurs, mais l'un d'eux était préfet, un autre 1\up{er} assistant et le troisième secrétaire. Que de soutiens enlevés à la fois à un édifice aussi fragile ! Il en fut sans doute ébranlé, mais grâce à Marie, notre bonne mère dont la protection ne nous a jamais manqué, la Congrégation fut assez heureuse pour résister à une crise aussi violente. MM. Jordan et Boni jugèrent convenable d'aller faire leurs études au séminaire de St Sulpice à Paris. Ils furent donc entièrement perdus pour nous. Heureusement M. Ripoud se décida à demeurer à Lyon. Il entra au grand séminaire qui était encore à la Providence pour y faire sa théologie. Les supérieurs eurent la condescendance de le laisser libre de continuer à administrer la Congrégation du fond de sa retraite. Aussitôt après le départ de nos deux autres amis il les remplaça dans leurs fonctions suivant le droit que lui en donnait le règlement ; les assistants se transportaient souvent au séminaire pour aller y prendre les ordres du préfet et tout continua à marcher comme auparavant. Il est sensible que
% Source : PDF original, vue 051, côté droit, page manuscrite 103.
cette position singulière où le préfet ne pouvait paraître que très rarement dans la Congrégation et se voyait forcé de lui faire parvenir ses ordres par l'intermédiaire des assistants dut être gênante. Dieu permit que bien au contraire elle fût fort avantageuse ; cette autorité dont les oracles étaient prononcés du fond de son sanctuaire n'en étaient que mieux respectés et mieux obéis. Les congréganistes n'eurent peut-être jamais plus d'attachement et de vénération pour leur préfet que pendant le temps qu'ils furent privés de sa présence.

Les vides que le départ de nos amis avait occasionnés furent promptement comblés par des réceptions nouvelles et la Congrégation se releva plus active et plus vivace que jamais.

\section*{Congrégation des hommes}
% Mention marginale de la page manuscrite 103.
Cependant le P. Roger tout en nous donnant tous ses soins ne négligeait rien pour faire fructifier le talent qui avait été remis entre ses mains. À peine eut-il fondé et donné la première impulsion à la Congrégation des demoiselles qu'il songea à un établissement pour les hommes mariés ; celle-ci prit naissance le jour de Noël 1802 sous le titre de Congrégation de la Ste Famille. Elle fut simplement divisée en deux branches ou sections ; la première s'appelait section du zèle. Je suis un peu embarrassé pour vous dire quelles étaient ses fonctions ; à vous dire la vérité je crois qu'elles se réduisaient à peu de choses. Quant à la seconde section sous le nom de section de charité elle soulageait un grand nombre de pauvres honteux.

En janvier 1803 le P. Roger compléta à peu près le système des Congrégations qui embrassa dès lors tous les
% Source : PDF original, vue 052, côté gauche, page manuscrite 104.
âges et tous les sexes en fondant la Congrégation des dames. Cette Congrégation fut divisée en sections dont les œuvres étaient analogues à celles de la Congrégation des jeunes gens. On leur donna les mêmes noms.

\section*{Congrégation des dames}
% Mention marginale de la page manuscrite 104.
D'après notre jeune âge et le peu de consistance que des jeunes gens dont la plupart n'avaient pas encore pris position dans le monde pouvaient avoir, il eût semblé qu'entre toutes ces Congrégations nous devions occuper le dernier rang et passer après tous les autres, mais il n'en était pas ainsi dans l'esprit du P. Roger. Nous étions ses aînés et il avait conçu pour nous une bien vive affection. La tendresse paternelle avait excité en lui à notre égard une confiance qu'on pourrait dire avoir été poussée au-delà des bornes où elle eût dû se maintenir. Il n'aurait pas établi une nouvelle Congrégation sans prendre l'avis et soumettre ses règlements aux chefs de la Congrégation des jeunes gens.

\section*{Opinion du P. Roger sur la Congrégation}
% Mention marginale de la page manuscrite 104.
Ce n'était pas chez lui le seul côté où on pouvait reconnaître le cœur d'un bon père qui se complaît à envisager ses enfants sous l'aspect le plus favorable. Il avait conçu une haute idée des destinées futures de la Congrégation et il ne négligeait rien pour nous faire partager à cet égard sa manière de voir. Il cherchait à bien pénétrer nos cœurs de la pensée que quelque foible qu'elle fût dans ses premiers moments elle pouvait un jour en prenant son essor rendre d'immenses services à la religion. Que de fois ne nous a-t-il pas répété : voyez, mes amis, ce petit noyau est jusqu'à ce jour bien peu de chose ; il est à peu près inaperçu ; mais patience, laissez agir la Providence ; Dieu a ses vues ; attendons que les moments marqués par le ciel soient arrivés et vous verrez
% Source : PDF original, vue 052, côté droit, page manuscrite 105.
alors ! Un jour viendra où vous serez étonnés vous-même de tout le bien qui se sera opéré par suite de la fondation de cette petite Congrégation. Lorsqu'en se rappelant ces paroles on reporte en même temps sa pensée sur certaines circonstances de l'histoire de la Congrégation dont nous aurons à nous occuper plus tard, on est bien tenté de croire qu'il avait plu au Seigneur de les faire connaître d'avance à notre guide.

C'est ainsi que notre vénérable père s'efforçait d'élever nos âmes et de faire naître en nous cet esprit de foi qui peut transporter les montagnes afin de nous rendre moins indignes d'accomplir plus tard malgré notre foiblesse tout le bien que la Providence pourrait demander de nous. Il voulait que nous fussions toujours disposés à nous laisser conduire par elle et que nous regardassions notre Congrégation comme devant être l'objet de tout notre respect et de toutes nos affections.

\section*{Intronisation de Mgr Fesch}
% Mention marginale de la page manuscrite 105.
Mgr Fesch prit au commencement de 1803 possession du siège archiépiscopal de Lyon. Il trouva grâce à l'administration du conseil qui avait été chargé par Mgr de Marbœuf de gouverner ce diocèse les affaires de la religion bien plus avancées qu'elles ne l'étaient partout ailleurs. Soixante théologiens réunis au grand séminaire vinrent assister à son intronisation et on comptait alors dans les petits séminaires déjà environ 300 jeunes gens qui se préparaient par les études littéraires aux études théologiques.

Le P. Roger fut parfaitement accueilli par le prélat qui même lui donna sa confiance et le choisit pour son confesseur. Les entrailles paternelles de ce bon père ne pouvaient pas oublier sa chère Congrégation ; il importait de la mettre en possession
% Source : PDF original, vue 053, côté gauche, page manuscrite 106.
de la bienveillance du premier pasteur du diocèse ; après avoir préparé les voies il présenta donc à son éminence le préfet et les assistants. Le cardinal les reçut avec bonté, parut prendre intérêt à la Congrégation et leur adressa des paroles d'encouragement.

\section*{Fin du schisme}
% Mention marginale de la page manuscrite 106.
Le premier soin du cardinal à son arrivée avait été de terminer l'organisation du diocèse. Dès que cette grande opération fut achevée et que tous les nouveaux pasteurs eurent pris possession de leurs postes il ne resta plus aucune apparence de schisme. Dès lors tout devait rentrer dans l'ordre normal ; il fut donc interdit de célébrer à l'avenir sans une permission spéciale les saints mystères hors des édifices consacrés au service divin. Rien n'était plus juste et plus nécessaire que cette mesure et pourtant elle affligea les fidèles ; on s'était fait une douce habitude de vivre pour ainsi dire en famille avec le bon Dieu.

Pour notre Congrégation en particulier cela élevait des obstacles à ses réunions religieuses qui nous firent craindre un instant de ne pas pouvoir continuer à célébrer en communauté nos solennités. Nous parvînmes heureusement à vaincre cette difficulté ; tant que d'autres circonstances n'ont pas rendu nos réunions impossibles nous avons toujours trouvé des chapelles particulières où on a eu la bonté de nous recevoir.

Les sœurs de St Vincent de Paul dites des Remparts d'Ainay ont été les premières chez lesquelles nous ayons reçu cette charitable hospitalité. Leur chapelle était fort petite, mais elle était alors bien suffisante pour notre petit nombre.
% Source : PDF original, vue 053, côté droit, page manuscrite 107.
C'était le jour de la fête de la Purification de la Ste Vierge 1803. Deux nouveaux amis venaient se joindre à nous. L'un d'eux est encore dans nos rangs ; mais je dois ici vous signaler la réception de l'autre comme un des faits mémorables de notre histoire car il est devenu une des gloires de la Congrégation et nous osons presque ajouter un de ses protecteurs dans le ciel.

\section*{Réception de Bertaud du Coin}
% Mention marginale de la page manuscrite 107.
Nommer Bertaud du Coin, Messieurs, c'est nommer un de ces chrétiens semblables à ceux des temps antiques animés comme eux d'un zèle ardent pour la défense de la foi. C'est nommer un de ces héros de la religion prêts à livrer leurs corps à toutes les tortures que pourraient inventer les tyrans plutôt que de s'écarter tant soit peu de la ligne du devoir. Jamais la fermeté de ce noble caractère ne s'est démentie. Dès l'âge de 16 ans il fut envoyé seul à Paris pour quelques affaires de famille et y demeura près d'un an libre de toute espèce de surveillance. Un séjour aussi périlleux pour un jeune homme encore dans l'adolescence ne produisit aucun mauvais effet sur du Coin. Sa piété fut toujours aussi vive et sa conduite aussi régulière. Il semble vraiment que le Seigneur voulut renouveler en sa faveur le miracle qu'il avait autrefois opéré en préservant de toute atteinte les trois jeunes hommes précipités dans une fournaise ardente. De même au milieu des flammes dangereuses attisées sans cesse par les passions et l'immoralité dans ce foyer de corruption dont notre saint ami était forcé de respirer l'air empesté sa belle âme conserva toute la fraîcheur de son innocence ; il s'abstint constamment de tout ce que la loi de Dieu interdit et notamment de la fréquentation des théâtres dans lesquels il n'a jamais mis les pieds. Il rentra donc dans
% Source : PDF original, vue 054, côté gauche, page manuscrite 108.
sa ville natale aussi pur qu'il en était parti. Sa vertu ne tarda pas à jeter un vif éclat. La Congrégation dut rechercher avec empressement la coopération d'un homme si bien fait pour la comprendre ; elle fut assez heureuse pour réussir. Pendant le petit nombre d'années que du Coin a pu partager les travaux de la Congrégation il a été certainement un des congréganistes qui lui a rendu le plus de services ; tour à tour chef de section ou assistant et enfin préfet, ainsi que nous aurons à le raconter plus tard, on ne sait ce que l'on doit le plus apprécier ou du zèle ou de l'excellent jugement dont il a constamment fait preuve. Nous vous signalerons bientôt sa conduite courageuse à l'époque où il partagea avec le souverain pontife les rigueurs de la persécution, mais comme l'occasion ne se représentera peut-être pas, je vais vous ajouter ici quelques mots sur l'époque où des devoirs d'état le tinrent éloigné de la Congrégation.

Au moment de la catastrophe des Cent-Jours lorsque du Coin eut acquis la conviction qu'il ne restait à Lyon aucun moyen de s'opposer à la conspiration qui replaçait Napoléon sur le trône, il se hâta de partir pour Paris ; là il s'enrôla dans ces volontaires qui s'armaient alors sous le nom de volontaires royaux pour la défense du roi légitime. On le chargea du commandement d'un détachement qu'il devait conduire à Gand. Sur la route il rencontre un régiment qui a déjà arboré la cocarde tricolore. Le régiment somme le détachement de crier vive l'empereur. Du Coin sans s'émouvoir répond par le cri français : Vive le roi ! Crie
% Source : PDF original, vue 054, côté droit, page manuscrite 109.
à sa troupe : en avant, et passe hardiment devant le front du régiment qui frappé de stupeur et d'étonnement pour cette noble audace semble pétrifié et demeure immobile.

Dieu a ses desseins ; quelquefois ils nous paraissent d'abord impénétrables ; on en reconnaît ensuite les raisons secrètes lorsqu'ils sont accomplis. La Providence avait destiné du Coin à devenir un apôtre pour la garde royale. D'officier de circonstance il devint après le retour du roi officier en titre dans cette garde. Bientôt il s'y fit admirer, on peut dire plus encore, il s'y fit vénérer, soit par son exactitude à remplir tous les devoirs de son nouvel état, soit par sa fermeté à repousser loin de lui tout ce que la loi de Dieu défend lors même que l'honneur si précieux aux militaires pouvait aux yeux de ses camarades s'y trouver compromis ; il sut habilement profiter de la haute considération qu'il avait acquise pour faire connaître, apprécier et goûter les vérités de la religion aux officiers dont il était entouré. Un très grand nombre lui durent le bonheur de revenir enfin à la pratique des devoirs religieux dont quelques-uns étaient éloignés depuis bien longtemps.

Les missionnaires de France donnèrent une mission aux régiments qui étaient en garnison à Versailles pendant le temps où notre ami était sous les drapeaux. Il seconda si bien leur zèle que ces apôtres de l'Évangile disaient hautement que dans cette circonstance le capitaine Bertaud (c'était le nom qu'on lui donnait alors) avait été plus missionnaire qu'eux tous et que c'était lui qui avait assuré le succès de la mission. Du Coin toujours zélé congréganiste et connaissant
% Source : PDF original, vue 055, côté gauche, page manuscrite 110.
tous les avantages qui pouvaient résulter d'une Congrégation voulut mettre le complément à ses travaux apostoliques en organisant dans la garde royale une Congrégation militaire. Je crois sans pouvoir cependant parfaitement l'assurer que cette Congrégation a subsisté jusqu'à la dissolution de la garde en 1830. Peu de temps après se trouvant en garnison à Rouen il essaya d'établir une Congrégation dans cette grande ville. J'ignore si ses efforts ont été couronnés d'un plein succès.

La carrière de notre saint confrère a été bien remplie, mais elle n'a pas été longue. Dieu acheva de le purifier par l'altération d'une santé qui jusque-là avait été très robuste. Il souffrit pendant plus d'un an avec une grande patience des douleurs de tête presque intolérables. Enfin une mort subite mais non pas imprévue vint le mettre en possession de ce beau ciel qui avait été l'objet constant de tous ses travaux comme il avait été celui de tous ses désirs.

Mme du Coin, la vénérable et sainte mère de l'ami dont nous venons de vous entretenir, connaissait tout l'attachement que son fils portait à la Congrégation ; elle désira après sa mort nous en donner une preuve. Elle voulut bien remettre comme souvenir à un membre de la Congrégation qu'il affectionnait particulièrement la bourse qui se trouva sur lui à ce dernier moment et destina à la Congrégation elle-même la petite somme que renfermait cette bourse. Le préfet jugea à propos pour perpétuer parmi nous la mémoire de ce congréganiste dans lequel la Providence nous a offert un modèle de consacrer cette somme à l'acquisition
% Source : PDF original, vue 055, côté droit, page manuscrite 111, avant le chapitre troisième.
d'une gravure représentant l'Assomption de la Ste Vierge. Cette gravure orne maintenant la salle de nos assemblées. En la considérant saluons Marie, remercions-la des grâces qu'elle a obtenues à notre ami et en nous rappelant les vertus dont il nous a donné l'exemple disons du fond du cœur : \textit{Sit memoria ejus in benedictione}\footnote{Traduction : « Que sa mémoire soit en bénédiction. » Référence : Siracide 46, 14 (formule apparentée, au pluriel dans la source biblique). (Note du transcripteur).}.
